\subsection{Modelo de Ciclo de Vida [Pendientes: 1. acordar duración y agenda de iteraciones, 2. definir flujo de revisión]}\label{sec:spmp-8-1}
El proyecto utilizará DAD con Inicio, Construcción y Transición, manteniendo la adaptación ya discutida por el equipo. Inicio aporta la delimitación, los requisitos y las primeras decisiones sobre arquitectura, datos y viabilidad, mientras que Construcción integra desarrollo de software y trabajo técnico mediante incrementos y Transición reúne la evaluación, las correcciones y la entrega. Esta secuencia no impide revisar decisiones previas cuando aparezca evidencia relevante. Por el contrario, las pruebas, la trazabilidad y la gestión de riesgos acompañarán el desarrollo, de manera que las incertidumbres sobre ejecución, datos o agregación puedan tratarse antes de comprometer la versión final~\cite{adaptacion2026}.

La figura \ref{fig:dad-pmp} resume la relación entre las fases sin asignarles duraciones proporcionales. Para el periodo de 2027-I, una iteración seleccionará un objetivo verificable, desarrollará el trabajo necesario, conservará sus comprobaciones y revisará el resultado para decidir la continuidad. Un incremento podrá producir una capacidad de la aplicación o evidencia que permita ajustar parámetros, datos o límites, por lo que el avance no se medirá únicamente por código escrito. Los productos y los criterios se identifican en las secciones \ref{sec:spmp-6-4} y \ref{sec:spmp-8-3}, mientras que los registros de ejecución seguirán las secciones \ref{sec:spmp-10-2} y \ref{sec:spmp-12-4}.
\begin{figure}[htbp]\centering
\begin{tikzpicture}[node distance=16pt,font=\normalsize,phase/.style={draw,rounded corners,text width=.84\linewidth,align=center,inner sep=7pt}]
\node[phase](a){\textbf{Inicio, 2026-II}\\Preparación y reducción de incertidumbres};
\node[phase,below=of a](b){\textbf{Construcción, 2027-I}\\Incrementos de software y evidencia técnica};
\node[phase,below=of b](c){\textbf{Transición, 2027-I}\\Evaluación, corrección y entrega};
\draw[-{Stealth}] (a)--(b);\draw[-{Stealth}](b)--(c);
\end{tikzpicture}
\caption{Ciclo de vida adoptado, sin escala temporal. Fuente: elaboración propia a partir de la metodología del equipo~\cite{adaptacion2026}.}\label{fig:dad-pmp}
\end{figure}

El procedimiento concreto de revisión de incrementos, sus participantes y su agenda formal siguen pendientes. La elección de DAD no establece por sí sola una asignación de responsabilidades ni autoriza a dar por acordado un flujo operativo específico.
\subsubsection{Análisis de Alternativas y Justificación - ¿Por qué pasamos de RUP a DAD? [Pendientes: 1. Completar comparación bibliográfica de ciclos y prácticas con las fuentes disponibles. 2. Definir y documentar las prácticas específicas de desarrollo que se utilizarán en el proyecto.]}\label{sec:spmp-8-1-1}
Rational Unified Process (RUP) fue considerado inicialmente por su naturaleza iterativa e incremental y su orientación a la arquitectura, características apropiadas para desarrollar la aplicación, construir la biblioteca de escenarios e integrar sus distintos componentes. Posteriormente se optó por DAD como una alternativa más reciente que conserva estas características pertinentes y permite adaptar la forma de trabajo a las necesidades del proyecto.

La elección de DAD responde a la necesidad de coordinar un producto de software con incertidumbres que no se resuelven únicamente mediante programación, puesto que la aplicación depende de decisiones sobre las entradas del simulador, su ejecución, el catálogo de escenarios y la aproximación espacial, por lo que revisar prioridades y obtener evidencia técnica desde las primeras iteraciones resulta coherente e importante para el trabajo previsto. La adaptación registrada conserva fases reconocibles y, al mismo tiempo, distribuye las pruebas, la arquitectura y la documentación entre las iteraciones; de este modo, la justificación de DAD frente a la alternativa inicialmente considerada se apoya en su adaptación al contexto y en la relación entre incertidumbre y resultados verificables~\cite{adaptacion2026}.

